\section{Return to Normalcy}

Over the weekend, someone uploaded child pornography to the Gornahoor forum. Bluehost disabled the entire site, without warning or explanation, until the situation could be evaluated by their legal department. They deleted the images and reactivated Gornahoor. When informed, I disabled attachments to forum posts. I also restricted posting, even in the public forum, to registered users. Even before I could complete that action, the offender managed to get in one last post, and in a taunting way wrote that he was glad the site was back up. He did leave a couple of links which I deleted.

On the plus side, it shows that there is one human activity that is almost universally considered reprehensible at a time when so much else is up for grabs. As for the forum, I'm inclined to abandon it. Its management is more trouble for very little benefit or use. From the logs, I can see dozens of attempts every day from China and Russia to hack into it; I have no conceivable notion of how that would benefit anyone.

In any event, I see that something needs to be done about web hosting. Gornahoor is getting too large for a mass market hosting service like Bluehost. On the other hand, there is not enough traffic for a virtual server. I have two months to find an alternative that is more reliable and faster.

In what follows, I want to make some quick points and mention future translations and reviews.

\paragraph{Salvation}


It did give me an opportunity to read the recent posts on Christianity, to which I'll provide some comments, ``dialectically'' of course. First of all, and most important, I have no intention of challenging anyone's faith. If anything you read here disturbs you or brings you doubt, I suggest you stop reading forthwith. This is decidedly an optional exercise and is not intended for everyone.

Nevertheless, if you need to persist and still feel agitated, I can make some calming points. Esoteric doctrine is not superadded to the exoteric, nor is it meant to be in conflict or opposition to it. It is our contention that theological (i.e., exoteric) dogmas can be restated in esoteric, or metaphysical, terms. The results speak for themselves: either you see it and agree with it or you don't.

I strongly recommend that everyone focus on his own salvation and not be so concerned about everyone else. Don't take your own for granted, actually don't even assume you even know what that means. We are not nominalists, so for us salvation requires a real change, an ``intellectual'' conversion, and not a nominal, or ``imputed'' one. There is no mechanical process involved and each man must follow the narrow way the best he can. If you are really so concerned about whether some fellow in Nepal is ``saved'' or not, then become a missionary.

\paragraph{Certainty}


What, then, is esoteric? When Augustine claims that truth is found in man's interiority, is he being esoteric? When St Gregory of Nyssa adopts Plato's and Aristotle's understanding of the soul, is he being esoteric? Where today will you even hear such teachings, outside of a survey course at university? The teaching is exoteric when taken in an external sense, as just one theory of the soul among many. The same teaching is esoteric when a man in his interiority actually ``sees'' the gradations and faculties of his own soul. In other words, esoterism is hidden in broad daylight.

That is why esoterism is not solely about ``thinking more carefully'' as Prof Cutsinger seems to believe. Actually, it is not about thinking at all, but rather about developing the powers of intellectual intuition. I don't know how this can be missed, since it is a fundamental teaching of every writer on Tradition. Nevertheless, it usually is.

\paragraph{The One Thing Needful}


Dialectic is weakness, we today need men of decision, men of conviction who are infallible in their understanding of metaphysics, intransigent and unsentimental in the face of truth, detached from personal concerns, and for whom tradition is not a career, nor raw material for journal articles, nor the road to money, sex, and power.

This is not for the faint or heart in our time, since deviation exists everywhere. What outwardly appear to be natural alliances are in reality false friends. The inner life is not matching the outer life. In a recent article on Takimag\footnote{\url{http://takimag.com/article/my\_otherkin\_headmate\_is\_a\_two\_spirited\_starseed\_kathy\_shaidle\#axzz2MjMZsZOX}}, cases of corporeal and spiritual mismatching are becoming common today. However, this was foreseen by Julius Evola 70 years ago, in this passage, the entirety of which will be sent to the subscribers to the S I mailing list:

\begin{quotex}
The castes, in which the ``race of the soul'' and the ``inner vocation'' do not correspond to the race of the body, so as those of every romantic laceration, in the last analysis, from the metaphysical point of view, are to be explained in such a way. Even modern psychology knows, at this point, so called ``second personalities''. And the more the minor forces diverge from the central direction, the more we will have, as the effect, men in whom the physical is not in accord with the soul, in which the spirit contrasts with the body and the soul, in which the vocation does not correspond to race or caste, the ``personality'' is in dissension with tradition, and so on.
\end{quotex}


Fortunately, Gornahoor has matured to reach a higher quality of reader, at least in my estimation of quality. This is obvious in the comments and posts, but especially in the private communications I have received. I hope that Goranhoor will continue to prove itself worthy of such readers.

\paragraph{Translations}


I will continue with Guido De Giorgio's description of the constitution of a traditional society. After the warriors, there will be the workers, and then the leader ``Capo'' . By way of background, De Giorgio uses the terms Silence, Rhythms ``ritmi'' , and Forms in a specific sense. They correspond to spirit, soul, and body, or to the concerns of the priests, warriors, and workers. Without going into a long digression, they can be thought of, without too much distortion, as the Tao, Duality, and the 10,000 things. The Silence is the ineffable at the center of the cross (in Guenon's sense), Rhythms are the opposing points (life and death, ascension and resurrection), and Forms are the manifested things.

I will also provide an article from the UR group called ``Plotinus, the greatest of pagan wisdom''. Please let me know if this is already available. Also some more Maurras, including his ``Latinity'' to complement De Giorgio's Romanity. S I subscribers will continue to receive the translation of Sintesi. Note that this is a direct translation from the Italian, and not a retranslation of the German translation.

\paragraph{Book Reviews}


I will also try some more book reviews, although the previous one simply faded away. They will be selected to elucidate certain points. These are the current candidates:

\begin{itemize}


\item \textit{St Gregory of Nyssa on the Human Soul}, by Constantine Cavarnos

\item \textit{Dante \& Aquinas}, by Philip Henry Wicksteed

\item \textit{Humanity, its Destiny and the Means to Attain It}, by Heinrich Denifle

\item Works of Mouni Sadhu
\end{itemize}



\flrightit{Posted on 2013-03-06 by Cologero}

\begin{center}* * *\end{center}

\begin{footnotesize}
\begin{sffamily}
\texttt{Sandra Eggers on 2013-03-06 at 00:39 said:}

Mouni Sadhu would be interesting.

\hfill

\texttt{Super Tutor on 2013-03-06 at 00:41 said:}

I wonder what will happen to the pervert?

\hfill

\texttt{Avery Morrow on 2013-03-06 at 03:26 said:}

Dialectic is weakness. Agreed.

I've noticed that a great majority of Evola readers, excluding the Nouvelle Droite ethno-nationalists (who are not my concern), are people who are of mixed ethnic heritage, or second or third-generation immigrants from non-Christian countries, or otherwise racially/sexually confused; in other words, people who can intuit some metaphysics, but cannot simply join the traditionalist Christian communities around them because they have within them some ``minor forces diverging from the central direction'' as Evola says.

The concerns of these people are worthy of an entire book, which I intend to write. Traditionalism is not something suitable for blogs, and people who read these blogs are to some extent lost in a thicket, although I credit this one with doing a good job of pointing towards the way out.

\hfill

\texttt{Omar Chowdhury on 2013-03-06 at 10:20 said:}

Intellectual intuition is missed because those who who overlook it are drenched in thought, all they know of is thought, conceiving of anything higher would be an admission of their own ignorant myopia.

\hfill

\texttt{Scardanelli on 2013-03-06 at 12:29 said:}

As far as translations are concerned, would anything by Cesare Della Riviera be possible? There isn't much available in English, and I'm curious as to what he might offer. I've only read of him through Evola, as well as Godwin's novel The Forbidden Book.

\hfill

\texttt{Logres on 2013-03-06 at 14:32 said:}

The implication of this teaching on esoteric/exoteric is that those who attack exoteric religions, by implication, are mounting an assault on secret doctrine and Tradition, as well.

\hfill

\texttt{Jason-Adam on 2013-03-06 at 15:11 said:}

Avery,

I have never thought about your point before that way but upon meditation I think you are correct -- aside from the Nouvelle Droite most Evolians today do seem to be very confused and lost individuals -- I notice a trend of repressed homosexuals, not to mention open homosexuals, also alot of people who confuse being a warrior with being a bully.

A problem that I think we all need to confront is this : most traditional Catholics of the present day are of the priest-type mind whereas warriors seem to be more drawn to neopaganism and the ND. We need to have more masculine, combative Catholics of the old crusader type and revive the warrior spirituality. I think it'd be an interesting personal challenge for a few men to try live according to the rules of the Templars.

\hfill

\texttt{David on 2013-03-06 at 18:20 said:}

@Logres : Rules of the Templars applied to a divine mission that was granted in the microcosm, macrocosm and inner world. It would be very difficult to find something like this in the modern that can unite those three. I agree with the idea, but to implement it is something else. It was a unique conjonction in christian history.

\hfill

\texttt{David on 2013-03-06 at 18:21 said:}

I'm sorry I meant @Jason-Adam,

\hfill

\texttt{Cologero on 2013-03-06 at 19:21 said:}

RE: confusing being a warrior with being a bully

The shudra are equally violent, but there is difference, a holy difference, between the shudra bully and the warrior. That will come out in the next segment which, I think, will be a little shocking, especially for those with contemporary sentimental ideas about spirituality. He may even topping Evola in this regard. Also, be sure to catch the allusion to Dante when he discusses dueling.

\hfill

\texttt{Alexander Maximilian on 2013-03-13 at 10:36 said:}

I highly concur. This is the one blog that maintains high integrity when it comes to the dispensation of Traditionalist content. Real Spiritual/ Metaphysical concerns beyond the mere political theories promulgated on other ``traditionalist'' sites that miss the entire point if initiation and dutiful service. I highly applaud the work done.

\hfill

\end{sffamily}
\end{footnotesize}

